\chapter{Supplementary Details}
\label{app:details}

\section{Reproduction Record}
An appendix can preserve details that are necessary for reproduction but would
interrupt the main discussion. For the illustrative study, this includes the
collection version, query identifiers, judgement files and the settings used
for each configuration. The record should connect these inputs to the result
files described in Chapter~\ref{sec:experiments}.

A concise checklist helps the reader locate the required material:
\begin{enumerate}
  \item Record the exact collection and query versions.
  \item Save preprocessing and retrieval settings for both configurations.
  \item Retain ranked outputs and paired query-level evaluation scores.
  \item Document exclusions and unexpected conditions.
  \item State the command or procedure used to reproduce each reported table.
\end{enumerate}

An appendix should remain readable on its own while making its relationship
to the main text clear. Cross-references are preferable to manually typed page
numbers because the document may change during revision. This appendix is
labelled with a letter, and its sections use the corresponding letter prefix.

% Demonstration-only break.
\newpage
\section{Checking the Compiled Document}
This continuation page provides one final example of the running header and
centred footer. The opening page of Appendix~\ref{app:details} should show its
lettered title and a page number, but no running header. Its entry in the table
of contents should point to that opening page rather than to this continuation.

When replacing the sample content, inspect the compiled document after major
structural changes. Long headings, large tables and figures close to a page
boundary can change pagination. Those changes are normal, but captions and
cross-references should remain associated with the correct content. Remove
the demonstration-only page breaks to allow the actual thesis to flow naturally.

The cover and abstract retain their existing unnumbered appearance. The table
of contents and all pages after it display page numbers. Page counting
continues from the start of the document, so the first visible number need not
be one. This example does not introduce a separate Roman-numeral sequence for
front matter or restart Arabic numbering at the first chapter.
